\section{统一口径对照实证}\label{sec:empirical}

前七章的比较建立在文献之上。文献的绩效数字彼此不可比（第~\ref{sec:audit}~节），因此需要一个变量受控的对照实验来检验核心论断。本章报告这样一组实验。

\subsection{实验设计}

\textbf{核心设计：唯一自由度是符号}。两策略共享以下全部设定，仅信号符号不同：

\begin{itemize}[leftmargin=2.2em]
  \item \textbf{资产池}：中国期货主力连续合约 45 个品种（经换月跳空守门前复权）；A 股宽基指数 7 个；A 股个股 5083 只（剔除上市不足与长期停牌）。
  \item \textbf{信号}：$s_t(L)=\sum_{k=0}^{L-1}\ln(1+r_{t-k})$，即过去 $L$ 日累计对数收益。趋势取 $+s$，反转取 $-s$。窗口 $L\in\{5,10,20,60,120,250\}$。
  \item \textbf{口径}：时序（TS）——按 $\mathrm{sign}(s_t)$ 乘波动率缩放后归一到总敞口 1；截面（XS）——按 $s_t$ 排序，多前列 20\%、空后列 20\%，两端等权，总敞口 1。
  \item \textbf{调仓}：每 5 个交易日调仓一次；信号使用至 $t$ 日收盘的信息，仓位自 $t+1$ 日生效；持有期内权重随收益漂移，换手按目标权重与实际权重之差计算。
  \item \textbf{成本}：期货单边 2 个基点；A 股指数 5 个基点（ETF/股指期货口径）；A 股个股单边 15 个基点（含佣金、印花税与滑点）。
  \item \textbf{样本}：2014-01-02 至 2026-09-30（3100 个交易日）。
\end{itemize}

\textbf{为什么这组设计能检验核心论断}。第~\ref{sec:unity}~节的论断是“两策略共享同一信号，唯一区别是符号”。若这个论断成立，那么在同一 $L$ 上，趋势与反转的换手应当\emph{完全相同}，成本拖累应当\emph{完全相同}，毛收益序列应当\emph{严格互为相反数}，净收益的差异应当完全由“信号符号 + 常数成本扣项”解释。这构成两条可证伪的验证门（对应第~\ref{sec:reproduce}~节的 G1 与 G2）。

\subsection{门禁检验：论断是否成立}

\begin{table}[htbp]
\centering\footnotesize
\caption{期货时序口径下两策略的对称性检验（G1/G2）}\label{tab:gate}
\begin{tabular}{@{}>{\raggedright\arraybackslash}p{1.6cm}>{\raggedright\arraybackslash}p{1.6cm}>{\raggedright\arraybackslash}p{2.0cm}>{\raggedright\arraybackslash}p{2.0cm}>{\raggedright\arraybackslash}p{2.0cm}>{\raggedright\arraybackslash}p{2.0cm}@{}}
\toprule
$L$ & 方向 & 年化换手 & 成本拖累 & 毛利夏普 & 净夏普 \\
\midrule
\multirow{2}{*}{5} & 趋势 & 52.41 & 1.06\% & $+0.40$ & $+0.20$ \\
                    & 反转 & 52.38 & 1.02\% & $-0.40$ & $-0.60$ \\
\midrule
\multirow{2}{*}{10} & 趋势 & 35.14 & 0.72\% & $+0.45$ & $+0.32$ \\
                     & 反转 & 35.12 & 0.68\% & $-0.44$ & $-0.58$ \\
\midrule
\multirow{2}{*}{60} & 趋势 & 16.14 & 0.33\% & $+0.26$ & $+0.21$ \\
                     & 反转 & 16.13 & 0.32\% & $-0.26$ & $-0.32$ \\
\midrule
\multirow{2}{*}{250} & 趋势 & 8.83 & 0.18\% & $+0.27$ & $+0.24$ \\
                      & 反转 & 8.85 & 0.17\% & $-0.27$ & $-0.30$ \\
\bottomrule
\end{tabular}
\end{table}

表~\ref{tab:gate}~给出三个结论。

\textbf{G1 与 G2 通过}。同一 $L$ 上，趋势与反转的年化换手几乎完全相等（$52.41$ 对 $52.38$，差异来自持有期内权重漂移的微小差异），成本拖累几乎完全相等（$1.06\%$ 对 $1.02\%$）；毛利夏普严格互为相反数（$+0.40$ 对 $-0.40$，$+0.45$ 对 $-0.44$，等等）。这确认了实现中没有引入非对称处理，“唯一区别是符号”在代码层面是成立的。

\textbf{成本的非对称作用是本实验最干净的发现}。以 $L=5$ 为例：毛利夏普为 $+0.40$ 与 $-0.40$（严格镜像），但净夏普为 $+0.20$ 与 $-0.60$——\emph{不是} $+0.20$ 与 $-0.20$。原因是成本是一个\emph{不区分方向的常数扣项}：它从两侧各减去约 $0.20$ 个夏普单位量级的收益（对应 1.04\% 的年化拖累 / 5.2\% 的年化波动 $\approx 0.20$）。好的一侧被削去一截，差的一侧被推向更差。\textbf{成本惩罚的是交易，而不是方向}；谁交易得更多，谁就被惩罚得更多——但当两策略交易量相同时，成本会放大两者的差距。

\textbf{短窗口的成本效率}。$L=5$ 时换手高达 52 倍/年，成本拖累 1.06\%，占毛利的一半以上；$L=250$ 时换手降至 8.83 倍/年，成本拖累 0.18\%，仅占毛利的 10\%。这直接印证了第~\ref{sec:duality}~节对偶五“换手与成本敏感度”的判断。

\subsection{期货：时序占优、截面呈 U 形}

\begin{table}[htbp]
\centering\footnotesize
\caption{中国期货：趋势与反转的夏普随窗口 $L$ 的变化}\label{tab:fut}
\begin{tabular}{@{}>{\raggedright\arraybackslash}p{1.0cm}>{\raggedright\arraybackslash}p{1.9cm}>{\raggedright\arraybackslash}p{1.9cm}>{\raggedright\arraybackslash}p{1.9cm}>{\raggedright\arraybackslash}p{1.9cm}>{\raggedright\arraybackslash}p{1.9cm}>{\raggedright\arraybackslash}p{1.9cm}@{}}
\toprule
\multirow{2}{*}{\textbf{$L$}} & \multicolumn{3}{c}{\textbf{时序口径（TS）}} & \multicolumn{3}{c}{\textbf{截面口径（XS）}} \\
\cmidrule(lr){2-4}\cmidrule(lr){5-7}
 & 趋势净 & 反转净 & 差值 & 趋势净 & 反转净 & 差值 \\
\midrule
5   & $+0.20$ & $-0.60$ & $+0.80$ & $-0.10$ & $-0.20$ & $+0.10$ \\
10  & $+0.32$ & $-0.58$ & $+0.90$ & $+0.11$ & $-0.31$ & $+0.42$ \\
20  & $+0.21$ & $-0.38$ & $+0.59$ & $-0.12$ & $-0.01$ & $-0.11$ \\
60  & $+0.21$ & $-0.32$ & $+0.53$ & $-0.82$ & $+0.72$ & $-1.54$ \\
120 & $+0.09$ & $-0.17$ & $+0.26$ & $-0.64$ & $+0.56$ & $-1.20$ \\
250 & $+0.24$ & $-0.30$ & $+0.54$ & $-0.32$ & $+0.25$ & $-0.57$ \\
\bottomrule
\end{tabular}
\end{table}

\textbf{时序口径：趋势全窗口占优，U 形未出现}。在 $L=5$ 到 $L=250$ 的全部六个窗口上，时间序列趋势策略的净夏普为正（$0.09\sim0.32$），时序反转策略的净夏普为负（$-0.17\sim-0.60$）。按美国市场的经验（短频反转、中频动量、长频反转），$L=5$ 与 $L=10$ 上本应出现反转占优，但中国期货市场并未出现。

这一结果值得单独强调：它\emph{不支持}“尺度决定归属”在中国期货时序口径上的适用性。可能的解释有三：一是中国商品期货的趋势持续性本身较强（产业资金与投机资金的行为特征）；二是 T+1 与涨跌停制度压制了短周期的反向交易；三是换手成本在短窗口上过高（$L=5$ 时 52 倍/年），足以淹没任何微弱的短期反转信号。本文无法区分这三种解释——它们是第~\ref{sec:future}~节的遗留课题 Q1。

\textbf{截面口径：U 形出现}。截面口径呈现了与直觉一致的结构：$L=10$ 时趋势占优（$+0.11$ 对 $-0.31$），$L=20$ 起反转追上，$L=60$ 与 $L=120$ 时反转显著占优（$+0.72$ 与 $+0.56$），$L=250$ 时优势收窄（$+0.25$）。趋势与反转的差值从 $+0.42$（$L=10$）单调变为 $-1.54$（$L=60$），再收窄至 $-0.57$（$L=250$）。

这是本文全部实验中\emph{唯一}出现 U 形结构的口径。它说明“尺度决定归属”不是普适定律，而是\emph{口径与市场联合决定}的结论：在有一定宽度横截面、且资产间存在共同因子的市场上，短窗口的横截面动量与中长窗口的横截面反转可以并存。

\textbf{反转侧在高 $L$ 上的经济含义}。$L=60$ 的截面反转净夏普 $0.72$、年化 $7.57\%$、最大回撤仅 $-18.74\%$，是本次实验中期货侧唯一兼具收益与稳健性的结果。它的含义是：在中国商品期货上买入过去三个月跌得最多的品种、卖出涨得最多的品种，在扣除 2 基点单边成本后仍有正的、显著的（NW $t=2.59$）收益。

\subsection{A 股：截面反转强，时序近零，短窗口被成本摧毁}

\begin{table}[htbp]
\centering\footnotesize
\caption{A 股：截面口径（5083 只个股）趋势与反转的对照}\label{tab:asu}
\begin{tabular}{@{}>{\raggedright\arraybackslash}p{1.0cm}>{\raggedright\arraybackslash}p{2.1cm}>{\raggedright\arraybackslash}p{2.1cm}>{\raggedright\arraybackslash}p{2.1cm}>{\raggedright\arraybackslash}p{2.1cm}>{\raggedright\arraybackslash}p{2.1cm}@{}}
\toprule
\textbf{$L$} & \textbf{趋势净夏普} & \textbf{反转毛利夏普} & \textbf{反转净夏普} & \textbf{反转换手} & \textbf{反转成本拖累} \\
\midrule
5   & $-2.49$ & $+1.08$ & $-0.35$ & 79.80 & 12.31\% \\
10  & $-2.11$ & $+1.10$ & $+0.10$ & 57.74 & 9.10\% \\
20  & $-2.15$ & $+1.46$ & $+0.76$ & 42.77 & 7.08\% \\
60  & $-1.93$ & $+1.46$ & $+0.99$ & 28.09 & 4.69\% \\
120 & $-1.63$ & $+1.24$ & $+0.86$ & 21.55 & 3.52\% \\
250 & $-1.17$ & $+0.86$ & $+0.56$ & 16.37 & 2.60\% \\
\bottomrule
\end{tabular}
\end{table}

\textbf{截面反转全面占优，且 U 形未出现}。A 股个股的截面动量策略在全部窗口上都是大幅负收益（净夏普 $-1.17\sim-2.49$，毛利夏普 $-0.88\sim-1.48$），其镜像的反转策略毛利夏普为 $+0.86\sim+1.46$。这与 A 股“追涨杀跌亏损、买跌卖涨获利”的市场共识一致，也是本次实验中最强的单一信号。

但注意：\emph{优势随 $L$ 的变化不是 U 形，而是倒 U 形或平台形}——反转的毛利夏普在 $L=20$ 与 $L=60$ 达到峰值（均为 $1.46$），在两端（$L=5$ 的 $1.08$、$L=250$ 的 $0.86$）下降。这与“短期反转最强”的微观结构直觉相反，说明 A 股的反转更可能是\emph{中期}现象——对应估值或情绪的回归，而非买卖价差反弹。

\textbf{短窗口被成本摧毁}。$L=5$ 时，反转的毛利夏普高达 $+1.08$，但换手 79.80 倍/年带来 12.31\% 的年化成本拖累，净夏普被压到 $-0.35$——\emph{信号完全被成本吞掉}。这是全文“成本敏感度”对偶最直接的证据：一个毛收益夏普大于 1 的信号，因换手过高而在扣除现实成本后变成负收益。随着 $L$ 增大，换手下降、成本拖累从 12.31\% 降到 2.60\%，反转的净夏普才转正并在 $L=60$ 达到 $0.99$。

这一结果对实践有直接含义：\textbf{在 A 股做统计套利型的截面反转，短窗口不可行，中期窗口（$L=20\sim120$）是唯一可行的区间}。这也解释了为什么 A 股的反转策略在实践中往往表现为“月度或季度调仓的低频价值/反转”，而不是高频统计套利。

\textbf{A 股指数时序：接近零收益，且不构成有效比较}。7 个宽基指数的时序结果（趋势净夏普 $-0.33\sim+0.17$，反转 $-0.27\sim+0.07$）全部在零附近，且最大回撤普遍在 $-50\%$ 至 $-79\%$ 之间。需要指出这组结果\emph{证据力很弱}：7 个 A 股宽基指数的相关系数高达 0.85 以上，其“组合”实质上是单一 A 股 beta 的择时，缺乏分散效应，且指数本身不可直接交易（需通过 ETF 或股指期货实现）。因此本文不把这组结果作为论据，仅作为对照列出。

\subsection{实证小结：论断的三个层次}

把四组结果合起来，可以对引言的论断给出分层的检验结论。

\textbf{第一层（信号统一）得到强支持}。期货时序口径下，趋势与反转的换手几乎完全相同（差异 $<0.1\%$）、成本拖累几乎完全相同、毛利夏普严格互为相反数。这从实证上确认了“两策略共享同一信号、唯一区别是符号”的构造性论断——它不是一个类比，而是一个可以逐项验证的事实。

\textbf{第二层（尺度决定归属）只在特定口径成立}。U 形结构只出现在期货截面口径上；在期货时序与 A 股截面上均未出现。因此这一论断应当修正为：\emph{“尺度决定归属”成立的条件是存在一个足够宽的横截面、且资产间存在共同因子}。在时序口径（单资产对自身历史）上，归属由市场特性决定，而非由尺度决定。

\textbf{第三层（成本与容量差异）得到强支持，且是实践中最重要的一层}。成本对短窗口策略的摧毁在两处得到直接印证：期货时序 $L=5$（成本占毛利 50\%）、A 股截面 $L=5$（毛利夏普 $1.08$ 被压到 $-0.35$）。而长窗口的成本占比降到 10\% 至 30\% 区间。这与第~\ref{sec:duality}~节对偶五以及第~\ref{sec:execution}~节的执行清单完全一致。

最后必须重复本章开头的定位：这组实证是\emph{变量受控的相对比较}，不是绝对收益水平的估计；由于未做多重检验校正（六个窗口 × 两方向 × 四组口径 = 48 个组合）、A 股侧做空不可实现、除权修复存在残余污染（第~\ref{sec:selfaudit}~节），本章的数字不应被当作“策略能赚多少”的答案。此外，本章全部指标基于日频收益的算术平均，未对自相关做夏普比率的标准误校正；Lo(2002) 指出在收益存在自相关时，夏普比率的样本估计方差会偏离独立同分布情形下的公式\cite{lo2002}，因此表~\ref{tab:gate}~中“成本使两侧差距不对称”的量化幅度应视为方向性判断，而非精确估计。
